%!TEX root = ../graduate-work.tex
\section{Problem Statement}

The primary limitation of contemporary verifiers is GPU memory consumption:
bound-propagation algorithms (IBP, CROWN, $\alpha$-CROWN) require weight
matrices and linear-relaxation coefficient matrices to be stored in their
entirety on a single accelerator, rendering verification of large models
infeasible.

\textbf{Goal of this work} is to adapt parallelism techniques developed for
training large language models to the problem of formal neural network
verification, and to experimentally characterise the trade-offs between
memory savings, bound tightness, and compatibility with complete verification.

\textbf{Objectives:}
\begin{enumerate}
    \item Analyse the applicability of LLM training parallelism strategies---
    Tensor Parallelism, Pipeline Parallelism, and FSDP---to bound-propagation
    algorithms and identify which are worth adapting.

    \item Implement \textbf{Tensor Parallelism} in a fork of
    \texttt{auto\_LiRPA}: sharding of weight matrices and linear-relaxation
    $A$-matrices across multiple GPUs, integration with ONNX models, soundness
    verification, and peak memory measurement.

    \item Implement \textbf{FSDP} for bound propagation---layer-wise
    AllGather/free of weights without modifying the computational
    logic---and achieve bitwise identity of bounds with the single-GPU baseline.

    \item Integrate \textbf{FSDP with complete verification}
    ($\beta$-CROWN + Branch-and-Bound); conduct a fair comparison by
    eliminating the auto-batch-enlargement artefact so that both modes
    process an identical workload.

    \item Extend FSDP sharding to \textbf{convolutional layers}
    (\texttt{BoundConv}) and verify correctness on ResNet models from
    the VNN-COMP benchmark.

    \item Run experiments on VNN-COMP benchmarks (MNIST-FC, ViT,
    CIFAR-100 ResNet) and determine under what conditions each method
    yields meaningful memory savings.
\end{enumerate}
